% !TEX root = ../main.tex
% Phụ lục: slide dự phòng cho phần hỏi đáp (phản biện + bảo vệ).
% Không nằm trong mục lục và không tính vào tổng số trang của phần nội dung.

\section{Phụ lục}

\begin{frame}[plain]
  \centering
  \vfill
  {\usebeamerfont{title}\usebeamercolor[fg]{title}\Large\textbf{PHỤ LỤC}}\\[0.25cm]
  {\normalsize Slide dự phòng cho phần hỏi đáp}
  \vfill
  \footnotesize
  \begin{tabular}{@{}ll@{}}
    \textbf{PL1} & Đánh giá trên chính website đã khám phá --- có hợp lệ không? \\
    \textbf{PL2} & Hướng đến ai: thứ tổng quát hóa là pipeline, không phải model \\
    \textbf{PL3} & Task có khả thi không, ngân sách bước, và ba lớp xác nhận hoàn thành \\
    \textbf{PL4} & Bốn câu hỏi nghiên cứu: số liệu và mức độ trả lời \\
    \textbf{PL5} & Chi phí chi tiết toàn chu trình \\
    \textbf{PL6} & Khám phá có định hướng: frontier hai tầng \\
    \textbf{PL7} & Nguyên nhân thất bại và chân dung hai mô hình \\
    \textbf{PL8} & Hai phép đo bổ trợ: thêm dữ liệu và đổi mô hình nền \\
  \end{tabular}
  \vfill
\end{frame}

%%% ---------------------------------------------------------------
\begin{frame}{PL1 --- Đánh giá trên chính website đã khám phá: có hợp lệ không?}
  \footnotesize
  \begin{columns}[t]
    \begin{column}{0.49\textwidth}
      \textbf{Cái \emph{không} xảy ra --- ranh giới train/test}
      \scriptsize
      \begin{itemize}\setlength{\itemsep}{2pt}
        \item Pipeline \kw{chưa từng nhìn thấy} 655 task của benchmark --- task huấn
              luyện do agent tự đề xuất từ giao diện.
        \item CSDL được \kw{khôi phục nguyên bản} trước mỗi đợt đo.
        \item Không trajectory thu thập nào được dùng lại lúc chấm.
      \end{itemize}
      \vspace{0.04cm}
      \textbf{Cái \emph{có} xảy ra --- và là chủ ý}
      \scriptsize
      \begin{itemize}\setlength{\itemsep}{2pt}
        \item Huấn luyện và đánh giá trên \emph{cùng} website --- đây là
              \kw{giao thức chuẩn} của cả dòng nghiên cứu: Go-Browse khám phá 20 URL
              mỗi domain của WebArena rồi chấm trên chính 812 task WebArena.
      \end{itemize}
    \end{column}
    \begin{column}{0.49\textwidth}
      \textbf{Vì sao không đo trên ``app khác cùng lĩnh vực''?}
      \scriptsize
      \begin{itemize}\setlength{\itemsep}{2pt}
        \item Đó là bài toán \emph{khác}: tổng quát hóa ngoài miền --- và số liệu cho
              thấy \kw{chưa ai giải được}. Trên Online-Mind2Web (web mở):
      \end{itemize}
      \vspace{-0.05cm}
      \centering
      \begin{tabular}{@{}L{3.1cm} R{1.3cm}@{}}
        \toprule
        \textbf{Mô hình} & \textbf{SR} \\
        \midrule
        NNetNav-7B & 4,00\% \\
        Go-Browse-7B (SOTA ${<}10$B) & 5,33\% \\
        GPT-4o-mini & 9,33\% \\
        \bottomrule
      \end{tabular}\par
      \vspace{0.08cm}
      \raggedright
      \scriptsize Hai hệ thống chuyên biệt hóa trong miền đều \emph{thua} GPT-4o-mini
      khi ra web mở.
    \end{column}
  \end{columns}
  \vspace{0.1cm}
  {\scriptsize\textbf{Phát biểu đúng phạm vi.} Khóa luận \kw{không} tuyên bố mô hình
  tổng quát sang website chưa từng thấy --- đã ghi rõ ở mục ``Hợp lệ ngoại vi''. Và nếu
  phản đối ``cùng app thì không hợp lệ'' đứng vững, nó bác bỏ luôn Go-Browse lẫn
  NNetNav --- toàn bộ đường cơ sở để so sánh.\par}
\end{frame}

%%% ---------------------------------------------------------------
\begin{frame}{PL2 --- Hướng đến ai, và thứ gì mới là thứ tổng quát}
  \footnotesize
  \begin{columns}[t]
    \begin{column}{0.5\textwidth}
      \textbf{Người dùng đích}
      \scriptsize
      \begin{itemize}\setlength{\itemsep}{2pt}
        \item Nhóm/doanh nghiệp vận hành \kw{một số website cố định} --- trang quản trị
              nội bộ, CMS, diễn đàn riêng --- muốn một agent chạy tại chỗ.
        \item Họ không cần agent ``biết mọi trang thương mại điện tử''. Họ cần agent
              \emph{thạo đúng trang của họ}, rẻ, không gửi dữ liệu ra ngoài.
      \end{itemize}
      \vspace{0.1cm}
      \begin{block}{Câu chốt}
        \scriptsize
        Thứ tổng quát hóa là \kw{pipeline}, không phải \kw{model}. Model là sản phẩm
        cho \emph{một} website; pipeline mới là đóng góp tái sử dụng được cho website
        tiếp theo.
      \end{block}
    \end{column}
    \begin{column}{0.48\textwidth}
      \textbf{Vì sao định vị này khả thi về kinh tế}
      \scriptsize
      \vspace{0.05cm}
      \begin{tabular}{@{}L{3.4cm} R{1.5cm}@{}}
        \toprule
        Chi phí API (4 website) & \$54,21 \\
        Chi phí GPU (fine-tune + đánh giá) & \$24,40 \\
        \midrule
        \textbf{Cả chu trình} & \textbf{\$78,61} \\
        \textbf{Trung bình mỗi website} & \kw{$<$ \$20} \\
        \bottomrule
      \end{tabular}\par
      \vspace{0.1cm}
      \raggedright
      Đối chiếu: AutoWebGLM cần \kw{20 người gán nhãn làm một tháng} cho tập
      $\approx$10.000 trajectory.
      \vspace{0.06cm}

      $\Rightarrow$ Dưới \$20 mỗi website và \kw{không cần người gán nhãn}, dựng model
      riêng cho \emph{từng} website hợp lý hơn hẳn việc đuổi theo một model tổng quát.
    \end{column}
  \end{columns}
\end{frame}

%%% ---------------------------------------------------------------
\begin{frame}{PL3 --- Task có khả thi không, đủ bước không, dựa vào đâu nói đã xong?}
  \scriptsize
  \begin{block}{1. Làm sao biết task do LLM đề xuất là hợp lý và làm được trên giao diện đó?}
    Bộ đề xuất là một \emph{agent}: đọc trực tiếp cây trợ năng của trang thật và có thể
    thao tác thử trước khi đề xuất --- không đề xuất từ trí nhớ. Quan trọng hơn:
    \kw{tính khả thi được chứng minh bằng thực thi, không bằng lời} --- chỉ task có ít
    nhất một lượt giải \emph{thành công và qua kiểm chứng} mới vào dataset.
    $\Rightarrow$ 637 task bị chính agent báo bất khả thi; 3981 lượt thử còn
    \kw{2117 hợp lệ (53,2\%)}.
  \end{block}
  \vspace{-0.1cm}
  \begin{block}{2. Vì sao ngân sách 15 bước (task cục bộ) và 25 bước (semantic task)?}
    Solver chiếm \kw{85,5\%} trong 20.808 lượt gọi LLM $\Rightarrow$ ngân sách bước
    \emph{chính là} ngân sách tiền; đây trước hết là quyết định chi phí. Nhưng dữ liệu
    cho thấy nó \kw{không bóp méo phép đo}: episode thành công có trung vị \kw{6 bước}
    (thất bại 13), và \kw{không episode nào} đạt được mục tiêu rồi mới bị cắt vì hết
    bước --- ngân sách đóng vai trò \emph{thước đo}, không phải \emph{hình phạt}.
  \end{block}
  \vspace{-0.1cm}
  \begin{block}{3. Dựa vào đâu để nói một task đã hoàn thành?}
    Ba lớp, mỗi lớp bắt loại lỗi lớp trước không thấy:
    \textbf{(i)} solver tự khai (\texttt{send\_msg\_to\_user}) hoặc vị từ tất định;
    \textbf{(ii)} \kw{LLM-as-Judge} soi lại trạng thái trang cuối $\rightarrow$ hạ cấp
    \kw{801 lượt} tự khai sai (27\%);
    \textbf{(iii)} cổng ``hành động hữu hiệu'' bằng mã $\rightarrow$ loại
    \kw{104 trajectory} tuyên bố xong mà không đổi gì trên trang --- nhóm lỗi mà judge
    \emph{về nguyên tắc} không thấy được, vì trang cuối không hề thay đổi.
    \vspace{0.04cm}

    \kw{Lưu ý}: đó là kiểm chứng lúc \emph{thu thập}. Con số 24,6\% được chấm bằng
    \kw{bộ chấm chính thức của WebArena}, không dùng judge của nhóm.
  \end{block}
\end{frame}

%%% ---------------------------------------------------------------
\begin{frame}{PL4 --- Bốn câu hỏi nghiên cứu: số liệu và mức độ trả lời}
  \scriptsize
  \setlength{\tabcolsep}{3pt}
  \renewcommand{\arraystretch}{1.15}
  \begin{tabular}{@{}L{0.9cm} L{4.5cm} L{4.9cm} L{1.5cm}@{}}
    \toprule
    \textbf{Câu hỏi} & \textbf{Số liệu} & \textbf{Kết quả} & \textbf{Trạng thái} \\
    \midrule
    \textbf{CH1} & 455 trang; 3981 lượt thử; 146 skill / 57 operator &
    1760 trajectory $\rightarrow$ 8554 mẫu; \$54,21 (0,0308\$/traj) & \kw{Đạt} \\
    \textbf{CH2} & 655 task, cùng prompt/ngân sách/bộ chấm; McNemar &
    13,3\% $\rightarrow$ 24,6\%; có ý nghĩa \emph{cả 4/4} website & \kw{Đạt} \\
    \textbf{CH3} & Tính lại các công trình trên đúng 4 website &
    Ngang Go-Browse-7B (2,0 $<$ CI $\pm$3,3); rẻ hơn 3,3$\times$ & \kw{Đạt} \\
    \textbf{CH4} & Phân loại lỗi; đo hiệu ứng bổ sung dữ liệu &
    85--96\% lỗi thực thi; gấp 3 dữ liệu chỉ $+$2,1 đpt & \kw{Đạt phần chẩn đoán} \\
    \bottomrule
  \end{tabular}

  \vspace{0.12cm}
  {\scriptsize\textbf{Phần thành thật phải nói kèm.} \textbf{O2} hoàn thành một phần:
  chỉ 12/142 semantic task đến từ ghép chuỗi operator, phần còn lại từ seed năng lực
  website. \textbf{O4} hoàn thành một phần: mỗi task mới đo \emph{một} lần, và
  \kw{ablation đồ thị so với khám phá phẳng --- phép đo duy nhất chứng minh trực tiếp
  giá trị của cấu trúc đồ thị --- vẫn chưa chạy}.\par}
\end{frame}

%%% ---------------------------------------------------------------
\begin{frame}{PL5 --- Chi phí chi tiết toàn chu trình}
  \footnotesize
  \begin{columns}[t]
    \begin{column}{0.5\textwidth}
      \textbf{Chi phí API (giai đoạn khám phá)}
      \scriptsize
      \vspace{0.05cm}
      \begin{tabular}{@{}L{2.6cm} R{1.5cm}@{}}
        \toprule
        \textbf{Nhà cung cấp} & \textbf{USD} \\
        \midrule
        OpenAI & 42,29 \\
        DeepSeek & 11,92 \\
        \midrule
        \textbf{Tổng} & \textbf{54,21} \\
        \bottomrule
      \end{tabular}\par
      \vspace{0.12cm}
      \textbf{Chi phí GPU (thuê theo giờ)}
      \vspace{0.05cm}
      \begin{tabular}{@{}L{2.6cm} R{1.5cm}@{}}
        \toprule
        \textbf{Giai đoạn} & \textbf{USD} \\
        \midrule
        Fine-tuning & 8,30 \\
        Đánh giá 655 task & 16,10 \\
        \midrule
        \textbf{Tổng} & \textbf{24,40} \\
        \bottomrule
      \end{tabular}
    \end{column}
    \begin{column}{0.48\textwidth}
      \textbf{Đơn giá và đối chiếu}
      \scriptsize
      \begin{itemize}\setlength{\itemsep}{3pt}
        \item \kw{0,0308 USD} mỗi trajectory thành công (54,21 / 1760).
        \item Go-Browse: 975,57 USD / 9504 trajectory $=$ 0,1026 USD
              $\Rightarrow$ khóa luận \kw{rẻ hơn 3,3 lần}.
        \item Synatra (tổng hợp \emph{ngoại tuyến} từ tài liệu, không hề chạm website
              thật): 0,031 USD/mẫu --- \emph{gần như trùng}, dù HSGE đã gồm cả chi phí
              điều khiển trình duyệt, đề xuất, giải và chấm lại task trên môi trường sống.
      \end{itemize}
      \vspace{0.08cm}
      \begin{block}{Điểm đáng chú ý không phải tổng chi phí, mà là \emph{vị trí} của nó}
        \scriptsize
        Rào cản tính toán không nằm ở huấn luyện hay đánh giá, mà ở \kw{khâu thu thập
        dữ liệu} --- vốn tiêu thời gian điều khiển trình duyệt và số lời gọi API chứ
        không tiêu GPU.
      \end{block}
    \end{column}
  \end{columns}
\end{frame}

%%% ---------------------------------------------------------------
\begin{frame}{PL6 --- Khám phá có định hướng: frontier hai tầng}
  \begin{columns}[c]
    \begin{column}{0.58\textwidth}
      \includegraphics[width=\linewidth,height=0.70\textheight,keepaspectratio]{fig_frontier.pdf}
    \end{column}
    \begin{column}{0.40\textwidth}
      \footnotesize
      \textbf{Cơ chế Frontier hai tầng}
      \begin{itemize}\setlength{\itemsep}{4pt}
        \item \textbf{Tier 0}: URL do solver ghé thăm thực tế $\rightarrow$ ưu tiên hàng đầu.
        \item \textbf{Tier 1}: URL trích từ DOM $\rightarrow$ giới hạn $\le 3$ trang/khuôn dạng (chống tràn frontier).
        \item \textbf{Lọc trùng cấu trúc}: mã băm nội dung loại \kw{71/455} node trùng lặp.
        \item \textbf{Solver Guard}: dừng cứng khi lặp hành động mà trạng thái không đổi.
      \end{itemize}
    \end{column}
  \end{columns}
\end{frame}

%%% ---------------------------------------------------------------
\begin{frame}{PL7 --- Nguyên nhân thất bại và chân dung hai mô hình}
  \footnotesize
  \begin{columns}[t]
    \begin{column}{0.52\textwidth}
      \textbf{Nguyên nhân chính của episode thất bại}
      \scriptsize
      \vspace{0.05cm}
      \setlength{\tabcolsep}{3pt}
      \begin{tabular}{@{}L{2.3cm} C{0.85cm} C{0.85cm} C{0.85cm} C{0.9cm}@{}}
        \toprule
        & \textbf{Shop} & \textbf{Reddit} & \textbf{GitLab} & \textbf{Admin} \\
        \midrule
        Lỗi hành động & 45\% & 34\% & 45\% & 39\% \\
        Kẹt vòng lặp & 34\% & 37\% & 31\% & 35\% \\
        Tự báo bất khả thi & 14\% & 23\% & 20\% & 25\% \\
        \bottomrule
      \end{tabular}\par
      \vspace{0.08cm}
      \raggedright
      \kw{85--96\%} episode thất bại có $\ge 1$ hành động bị trình duyệt từ chối
      (2315 lỗi / 655 episode) --- phần lớn là quá hạn chờ thao tác trên phần tử
      \emph{có} hiển thị trong cây trợ năng.
    \end{column}
    \begin{column}{0.46\textwidth}
      \textbf{Fine-tune làm mô hình \emph{bền bỉ} hơn, không \emph{chính xác} hơn}
      \scriptsize
      \vspace{0.05cm}
      \begin{tabular}{@{}L{2.7cm} C{1.05cm} C{1.05cm}@{}}
        \toprule
        & \textbf{Nền} & \textbf{Sau FT} \\
        \midrule
        Phiên hỏng định dạng & 26,0\% & \kw{0,6\%} \\
        Tổng số bước & 5424 & 6727 \\
        Lỗi hành động / bước & 0,358 & 0,344 \\
        Trung vị bước (đỗ) & 4 & 6 \\
        Trung vị bước (trượt) & 8 & 13 \\
        \bottomrule
      \end{tabular}\par
      \vspace{0.08cm}
      \raggedright
      Tỷ lệ lỗi trên mỗi bước gần như \emph{không đổi}: mô hình mắc lỗi thao tác với
      cùng tần suất, chỉ là nó \kw{sống sót đủ lâu} để thao tác nhiều hơn.
      $\Rightarrow$ điểm nghẽn còn lại nằm ở \kw{tầng thực thi}.
    \end{column}
  \end{columns}
\end{frame}

%%% ---------------------------------------------------------------
\begin{frame}{PL8 --- Hai phép đo bổ trợ: thêm dữ liệu và đổi mô hình nền}
  \footnotesize
  \begin{columns}[t]
    \begin{column}{0.5\textwidth}
      \textbf{(a) Thêm dữ liệu thì được gì?}
      \scriptsize
      \vspace{0.05cm}
      \setlength{\tabcolsep}{3pt}
      \begin{tabular}{@{}L{1.9cm} C{1.1cm} C{1.0cm} C{1.0cm}@{}}
        \toprule
        \textbf{Website} & \textbf{Mẫu thêm} & \textbf{Trước} & \textbf{Sau} \\
        \midrule
        Reddit & $+$1856 & 20,8\% & 23,6\% \\
        Shopping Admin & $+$2410 & 15,4\% & 17,0\% \\
        \midrule
        \textbf{Cộng (288 task)} & \textbf{$+$4266} & \textbf{17,4\%} & \textbf{19,4\%} \\
        \bottomrule
      \end{tabular}\par
      \vspace{0.08cm}
      \raggedright
      Thêm dữ liệu \emph{có} nâng điểm, nhưng hiệu suất thấp và giảm nhanh: gần
      \kw{gấp ba} lượng dữ liệu chỉ đổi lấy \kw{$+$2,1 đpt}. Website nhận
      \emph{nhiều} mẫu hơn lại đổi được \emph{ít} điểm hơn trên mỗi mẫu.
    \end{column}
    \begin{column}{0.48\textwidth}
      \textbf{(b) Đổi mô hình nền, giữ nguyên tập dữ liệu}
      \scriptsize
      \begin{itemize}\setlength{\itemsep}{2.5pt}
        \item Cùng 8554 mẫu, fine-tune Qwen-2.5-7B (đúng mô hình nền Go-Browse dùng):
              \kw{83/655 $=$ 12,7\%} so với 8,3\% của chính nó khi chưa fine-tune.
        \item \kw{Không dùng con số này để so cùng mô hình nền với Go-Browse}: toàn bộ
              275 phản hồi ``hỏng định dạng'' đều là lệnh \texttt{report\_infeasible}
              \emph{viết trần} (92--98\%), bị bộ phân giải loại bỏ --- bản 7B chỉ ghi
              nhận 10/182 lần báo bất khả thi ở Shopping Admin so với 41 của bản 9B.
        \item $\Rightarrow$ 12,7\% là \kw{cận dưới do công cụ đo đặt ra}. Điều rút ra
              được: tập dữ liệu nâng được \emph{cả hai} mô hình nền, nên nó không gắn
              riêng với một kiến trúc.
      \end{itemize}
    \end{column}
  \end{columns}
\end{frame}
